\documentclass[11pt, a4paper]{article}

% Packages
\usepackage[utf8]{inputenc}
\usepackage{geometry}
\geometry{top=1in, bottom=1in, left=1in, right=1in}
\usepackage{amsmath, amssymb, amsfonts, bm}
\usepackage{hyperref}
\usepackage{graphicx}
\usepackage{booktabs}
\usepackage{titlesec}
\usepackage{mathpazo} % Palatino font for a professional look

% Hyperref setup
\hypersetup{
    colorlinks=true,
    linkcolor=blue,
    filecolor=magenta,      
    urlcolor=cyan,
    pdftitle={Jointly Optimized Latent Bottleneck for Edge-Efficient Listwise Reranking},
}

\title{\textbf{Jointly Optimized Latent Bottleneck for Edge-Efficient Listwise Reranking}}
\author{Phạm Huy Tùng}
\date{\today}

\begin{document}

\maketitle

\begin{abstract}
Deploying high-precision listwise reranking on edge devices presents a fundamental conflict between retrieval accuracy and hardware constraints. While Late-Interaction architectures achieve state-of-the-art recall by preserving fine-grained token representations, their $\mathcal{O}(N)$ memory scaling makes them unviable for local indexing. Previous attempts to compress these representations introduce critical trade-offs: spatial chunking fractures global reasoning, while pure intent extraction prunes unqueried, high-entropy lexical facts necessary for listwise tie-breaking. This proposal introduces a unified, edge-efficient Late-Interaction framework that employs Entropy-Driven Compression to dynamically bound documents into an adaptive $\mathcal{O}(256)$ vector budget via a gated Cross-Attention bottleneck. By training a Residual Query Adapter and the Document Compressor jointly via a Two-Teacher distillation pipeline (a generative LLM for complex queries and a Cross-Encoder for target scoring), and introducing an IDF-weighted Chamfer divergence, the architecture guarantees the preservation of both global multi-hop intent and microscopic lexical grounding without manifold drift.
\end{abstract}

\section{Introduction and Motivation}
Deploying high-precision listwise reranking on edge devices requires bridging the gap between state-of-the-art retrieval accuracy and strict local hardware constraints. While architectures utilizing the MaxSim Late-Interaction operator (e.g., ColBERT) excel at listwise comparison by maintaining token-level fidelity, their memory footprints grow linearly with document length. 

Current compression paradigms fail to support the granular demands of listwise reranking. Spatial, paragraph-level chunking preserves exact lexical facts but fractures global logic and coreferences. Conversely, query-driven intent extraction successfully resolves multi-hop reasoning but aggressively prunes unqueried parameters, dates, and rare nouns---the exact microscopic evidence required to sort highly similar documents. We propose a jointly optimized Latent Bottleneck that forces a neural network to bypass this trade-off natively, carving out an adaptive vector budget that dynamically acts as both a reasoning engine and a strict spatial dictionary.

\section{Architectural Framework}
The framework operates on a dual-encoder Late-Interaction paradigm, modified to enforce aggressive, length-agnostic compression on the document side while maintaining exact-match fidelity on the query side.

\subsection{The Residual Query Adapter}
To cap edge Neural Processing Unit (NPU) latency, the user's short search query $Q$ is padded with \texttt{[MASK]} tokens up to a strict length of 32. It is processed by an unfrozen Query Adapter encoder $E_Q$. To prevent semantic drift and protect the model's ability to execute exact keyword matches, the adapter utilizes a residual skip connection. The non-linear transformations $F_Q$ refine the base embeddings $Q_{\text{emb}}$, which are then added back to the original tokens:
\begin{align}
    Q_{\text{emb}} &= \text{Embedding}(\text{PadWithMask}(Q, 32)) \in \mathbb{R}^{32 \times d} \\
    Q_{\text{latent}} &= \text{LayerNorm}\left( Q_{\text{emb}} + F_Q(Q_{\text{emb}}) \right) = \{q_1, \dots, q_{32}\} \in \mathbb{R}^{32 \times d}
\end{align}

\subsection{The Latent Bottleneck Document Compressor}
To bypass the $\mathcal{O}(N)$ memory explosion, the compressor decouples vector allocation from literal text length through three distinct phases:
\begin{enumerate}
    \item \textbf{Contextualization:} The document $D$ is passed through a deep Transformer backbone $E_D$, utilizing stacked self-attention layers to produce full, uncompressed contextualized token embeddings:
    \begin{equation}
        H_D = E_D(D) = \{h_1, \dots, h_{L_D}\} \in \mathbb{R}^{L_D \times d}
    \end{equation}
    
    \item \textbf{Cross-Attention Extraction:} A fixed matrix of $M=256$ learnable summary tokens $S = \{s_1, \dots, s_{256}\}$ acts as query probes in a Multi-Head Cross-Attention (MHCA) block $C$. This extracts semantic summaries, seamlessly resolving global coreferences across the contextualized space $H_D$:
    \begin{equation}
        B_D = C(S, H_D) = \{b_1, \dots, b_{256}\} \in \mathbb{R}^{256 \times d}
    \end{equation}
    
    \item \textbf{Entropy-Driven Adaptive Gating (Soft-Pruning):} A linear gating head evaluates the entropy of the bottleneck matrix, predicting a continuous retention probability $\pi_k$ for each vector. Soft-masking scales the vectors during training to maintain differentiability, allowing the active capacity to shrink below the 256 upper bound based on document density:
    \begin{align}
        \pi_k &= \sigma(\mathbf{w}_g^\top b_k + w_{bias}) \in (0, 1) \\
        \tilde{B}_D &= \pi_D \odot B_D = \{\pi_1 b_1, \dots, \pi_{256} b_{256}\}
    \end{align}
\end{enumerate}

\section{Interaction and Scoring}
The model evaluates relevance using an asymmetric Late-Interaction (MaxSim) operator. The 32 query vectors scan the soft-masked bottleneck vectors, calculating the maximum dot-product for each query token. This scoring mechanism is computationally bounded by $32 \times 256$ operations per document, ensuring deterministic sub-millisecond edge latency:
\begin{equation}
    S(Q, D) = \sum_{i=1}^{32} \max_{b_j \in \tilde{B}_D} \left( q_i^\top b_j \right)
\end{equation}

\section{Two-Teacher Distillation and Joint Loss Formulation}
The architecture relies on a Two-Teacher Distillation paradigm to decouple complex query generation from discriminative scoring. 

First, a \textbf{Generative Teacher} (a 70B parameter LLM) reads the corpus to perform Inquisitive Generation, outputting a synthetic dataset containing both direct (fact-seeking) and indirect (multi-hop relational) queries. Second, a \textbf{Discriminative Teacher} (a high-capacity Cross-Encoder) evaluates these query-document pairs, concatenating them through its full self-attention layers to establish highly calibrated, continuous target margins ($\Delta s^*$). 

The edge-efficient Student pipeline is trained jointly on this setup. The total loss enforces four distinct behaviors, dynamically allocating the vector budget:

\subsection{Query-Driven Intent Extraction (Margin-MSE)}
Forces the bottleneck vectors to encode global logic, answer synthetic queries, and explicitly mimic the target margins predicted by the Cross-Encoder Teacher:
\begin{align}
    \Delta s &= S(Q, D^+) - S(Q, D^-) \\
    \mathcal{L}_{\text{Margin}} &= (\Delta s - \Delta s^*)^2
\end{align}

\subsection{Spatial and Lexical Grounding (IDF-Weighted Chamfer)}
Acts as a geometric tether against the "Teacher Blind Spot." It calculates the Cosine distance strictly between the compressed vectors $\tilde{B}_D$ and the pre-compression contextualized tokens $H_D$. Cosine similarity prevents manifold drift, forcing the bottleneck to physically memorize high-entropy lexical details required for listwise reranking:
\begin{align}
    \mathcal{L}_{\text{coverage}} &= \frac{1}{L_D} \sum_{j=1}^{L_D} \text{IDF}(w_j) \cdot \min_{k \in \{1, \dots, 256\}} \left( 1 - \frac{h_j^\top \tilde{b}_k}{\|h_j\| \|\tilde{b}_k\|} \right) \\
    \mathcal{L}_{\text{precision}} &= \frac{1}{\sum \pi_i} \sum_{k=1}^{256} \pi_k \cdot \min_{j \in \{1, \dots, L_D\}} \left( 1 - \frac{\tilde{b}_k^\top h_j}{\|\tilde{b}_k\| \|h_j\|} \right) \\
    \mathcal{L}_{\text{Chamfer}} &= \mathcal{L}_{\text{coverage}} + \beta \mathcal{L}_{\text{precision}}
\end{align}

\subsection{Semantic Consistency Loss (Query Alignment)}
Guarantees that direct and indirect queries targeting the same underlying factual anchor do not fracture the latent space. It enforces an MSE penalty if the MaxSim scores for a direct and indirect query pair diverge when evaluating the same positive document $D^+$:
\begin{equation}
    \mathcal{L}_{\text{Consistency}} = \left( S(Q_{\text{direct}}, D^+) - S(Q_{\text{indirect}}, D^+) \right)^2
\end{equation}

\subsection{Information Entropy Penalty (Adaptive Sparsity Budget)}
An $L_1$ regularization applied exclusively to the gating probabilities ($\pi_k$), which allows the architecture to execute Entropy-Driven Compression and generalize across document lengths. For short, sparse documents (e.g., MS MARCO), the penalty successfully pushes the majority of $\pi_k \to 0$. For dense, lengthy documents (e.g., Wikipedia), the reconstruction gradients from $\mathcal{L}_{\text{Chamfer}}$ and $\mathcal{L}_{\text{Margin}}$ naturally overpower the sparsity penalty, keeping more gates open to capture necessary facts without exceeding the hardware ceiling:
\begin{equation}
    \mathcal{L}_{\text{Sparsity}} = \sum_{k=1}^{256} \pi_{D^+}^{(k)} + \sum_{k=1}^{256} \pi_{D^-}^{(k)}
\end{equation}

\subsection{Final Total Objective}
The pipeline is optimized jointly by minimizing the sum of the four objectives:
\begin{equation}
    \mathcal{L}_{\text{total}} = \mathcal{L}_{\text{Margin}} + \lambda_1 \mathcal{L}_{\text{Chamfer}} + \lambda_2 \mathcal{L}_{\text{Consistency}} + \lambda_3 \mathcal{L}_{\text{Sparsity}}
\end{equation}

\end{document}